\documentclass[10.5pt]{ctexart}
\usepackage[a4paper,top=1.45cm,bottom=1.45cm,left=1.55cm,right=1.55cm]{geometry}
\usepackage{hyperref}
\usepackage{enumitem}
\usepackage{multicol}
\usepackage{titlesec}
\usepackage{xcolor}
\usepackage{fancyhdr}
\setlength{\parindent}{0pt}
\setlength{\parskip}{4pt}
\setlist{nosep,leftmargin=1.3em}
\definecolor{ormuted}{HTML}{626262}
\definecolor{orline}{HTML}{B9B5AE}
\hypersetup{colorlinks=true,urlcolor=black,linkcolor=black}
\pagestyle{fancy}
\fancyhf{}
\lhead{\small Opportunity Radar · OR Morning · 2026-10-19}
\rhead{\small\thepage/8}
\cfoot{}
\renewcommand{\headrulewidth}{0.35pt}
\titleformat{\section}{\Huge\bfseries}{ }{0pt}{}
\begin{document}
{\small\bfseries OR MORNING / SECTION 01 / 2026-10-19}
\vspace{2mm}
\hrule height 0.9pt
\vspace{1.5mm}
\hrule height 0.25pt
\vspace{4mm}
\section*{今日总览}
{\color{ormuted} 如果今天只能看三分钟，这一版必须已经足够。}
\vspace{2mm}
{\fontsize{25}{29}\selectfont\bfseries 在职课程培训班}\par
\vspace{2mm}
{\large 本版筛出 7 个高信号事件。}
\vspace{3mm}
{\Large\bfseries 在职课程培训班}\par
首页 学院概况 学院简介 数说学院 现任领导 学术委员会 行政机构 新闻公告 学院新闻 学院公告 师生风采 光影视界 师资队伍 专任教师 博士后队伍 人才招聘 人才培养 本科生 学术型硕士研究生 专业学位硕士研究生 学术型博士研究生 专业学位博士研究生 在职课程培训班 同等学力申请硕士学位 科学研究 研究机构 研究成果 合作交流 对外合作 国际交流 校友名录 光影记录 校友名录 高礼研究院 EN 教师入口 人才培养 本科生 学术型硕士研究生 专业学位硕士研究生 学术型博士研究生 专业学位博士研究生 在职课程培训班 同等学力申请硕士学位 在职课程培训班 您当前所在的位置： 首页 > 人才培养 > 在职课程培训班 19 2026-10 源自人大高瓴，成就AI英才 2026年（秋）人工智能在职课程培训班招生简章 19 2026-10 人大高瓴人工智能学院人工智能专业在职课程培训的10个常见咨询问答 01 2026-03 2026年中国
\textbf{与你有关：} 这是当前用户可实际参与或申请的高可触达机会；参与范围：北京 / 人大；命中用户订阅版面：opportunity、research、ai；来源接近一手且权威度高；标题/正文包含高信号变化
\textbf{建议：} 检查报名资格、截止时间、参与形式和预期收益；合适则立即加入任务系统。
{\footnotesize\color{ormuted} 来源：中国人民大学高瓴人工智能学院公告 · \url{https://ai.ruc.edu.cn/rcpy/zzkcpxb/index.htm}}
\vspace{3mm}
\hrule height 0.25pt
\vspace{3mm}
\begin{multicols}{2}
{\large\bfseries 姚班培养-清华大学交叉信息研究院}\par
首页 热点新闻 讲座预告 院系概况 院系简介 院长寄语 现任领导 院系沿革 新闻动态 院系动态 媒体聚焦 人员队伍 全职教师 研究系列 荣誉客座 兼职教师 博士后 行政人员 实验员 姚班本科生 姚班培养 姚班招生 研究生 研究生培养 研究生招生 量子中心 中心简介 师资力量 科研概况 量子中心动态 科学研究 组织架构 课题组介绍 研究动态 科研简报 校园生活 学生活动 衣食住行 公开招聘 教师招聘 博士后招聘 职员招聘 校友频道 校友风采 校友俱乐部 校友捐赠 姚班本科生 姚班培养 姚班招生 清华学堂计算机科学实验班 姚班介绍 “清华学堂计算机科学实验班（姚班）”由世界著名计算机科学家姚期智院士于2005年创办，是清华大学培养国际顶尖计算机科学及交叉创新人才的本科生实验班，被誉为“最优秀的本科生和最优秀的本科教育”标杆。 姚班面向科技发展前沿与国家重大需求构建培养体系，以计算机科学、人工智能和量子信息三大前沿专业方向为核心，供学
\textbf{与你有关：} 这是当前用户可实际参与或申请的高可触达机会；参与范围：北京 / 清华；命中用户订阅版面：opportunity、research、ai；来源接近一手且权威度高；标题/正文包含高信号变化
\textbf{建议：} 检查报名资格、截止时间、参与形式和预期收益；合适则立即加入任务系统。
{\footnotesize\color{ormuted} 来源：清华交叉信息研究院 · \url{https://iiis.tsinghua.edu.cn/ybbks/ybpy.htm}}
\vspace{2mm}
\hrule height 0.2pt
\vspace{2mm}
{\large\bfseries Back to Careers}\par
Back to Careers Join our team All Teams AI Research \& Engineering Applied AI Communications Compute Data Science \& Analytics Engineering \& Design - Product Finance Hardware Legal Marketing \& Brand People Product Management, Support, \& Operations Public Benefit Public Policy Safeguards (Trust \& Safety) Sales Security Software Engineering - Infrastructure Technical Education Technical Program Management All Offices Ban
\textbf{与你有关：} 参与范围：全球 / 多为海外全职；命中用户订阅版面：opportunity、ai、research；来源接近一手且权威度高；标题/正文包含高信号变化
\textbf{建议：} 检查报名资格、截止时间、参与形式和预期收益；合适则立即加入任务系统。
{\footnotesize\color{ormuted} 来源：Anthropic Careers · \url{https://www.anthropic.com/careers}}
\vspace{2mm}
\hrule height 0.2pt
\vspace{2mm}
{\large\bfseries Newsroom \textbackslash{} Anthropic}\par
Newsroom Press inquiries press@anthropic.com Non-media inquiries How to get support Media assets Download press kit Expanding the Cyber Verification Program Announcements Oct 6, 2026 We’re launching a new, expanded version of our Cyber Verification Program, which makes advanced cyber capabilities and reduced blocking classifiers available to qualifying security professionals. Announcements Oct 7, 2026 Introducing Cla
\textbf{与你有关：} 命中用户订阅版面：ai、research、math；来源接近一手且权威度高；标题/正文包含高信号变化
{\footnotesize\color{ormuted} 来源：Anthropic Research · \url{https://www.anthropic.com/news}}
\vspace{2mm}
\hrule height 0.2pt
\vspace{2mm}
{\large\bfseries Algorithms \& theory}\par
Algorithms and theory Google’s mission presents many exciting algorithmic and optimization challenges across different product areas including Search, Ads, Social, and Google Infrastructure. These include optimizing internal systems such as scheduling the machines that power the numerous computations done each day, as well as optimizations that affect core products and users, from online allocation of ads to page-vie
\textbf{与你有关：} 命中用户订阅版面：ai、research、math；来源接近一手且权威度高；标题/正文包含高信号变化
{\footnotesize\color{ormuted} 来源：Google Research Blog · \url{https://research.google/research-areas/algorithms-and-theory}}
\vspace{2mm}
\hrule height 0.2pt
\vspace{2mm}
{\large\bfseries v4.36.0-rc1}\par
v4.36.0-rc1 Pre-release Pre-release Compare Choose a tag to compare Sorry, something went wrong. Filter Loading Sorry, something went wrong. Uh oh! There was an error while loading. Please reload this page . No results found View all tags github-actions released this 16 Oct 19:43 · 77 commits to master since this release v4.36.0-rc1 625d891 chore: bump toolchain to v4.36.0-rc1 (\#44879) Co-authored-by: Anne C.A. Baane
\textbf{与你有关：} 命中用户订阅版面：math、engineering、research；来源接近一手且权威度高；标题/正文包含高信号变化
{\footnotesize\color{ormuted} 来源：mathlib Releases · \url{https://github.com/leanprover-community/mathlib4/releases/tag/v4.36.0-rc1}}
\vspace{2mm}
\hrule height 0.2pt
\vspace{2mm}
{\large\bfseries Offloaded inference for real-world physical AI robotics}\par
At a glance Challenges a core assumption in robotics AI: Our research shows that running physical AI inference exclusively on onboard GPUs can limit robot performance, battery life, and scalability, and that offloading inference to edge or cloud GPUs can offer significant advantages. Demonstrates measurable benefits of inference offloading: Across representative mobile manipulation workloads, offloading improved task
\textbf{与你有关：} 命中用户订阅版面：ai、research、engineering；来源接近一手且权威度高；标题/正文包含高信号变化
{\footnotesize\color{ormuted} 来源：Microsoft Research Blog · \url{https://www.microsoft.com/en-us/research/blog/offloaded-inference-for-real-world-physical-ai-robotics/}}
\end{multicols}
\newpage
{\small\bfseries OR MORNING / SECTION 02 / 2026-10-19}
\vspace{2mm}
\hrule height 0.9pt
\vspace{1.5mm}
\hrule height 0.25pt
\vspace{4mm}
\section*{机会雷达}
{\color{ormuted} 北京、全国与线上真正可报名的 Hackathon、竞赛、科研招募、实习、Workshop 与截止。}
\vspace{2mm}
{\LARGE\bfseries 在职课程培训班}\par
\vspace{1mm}
{\color{ormuted} 本版筛出 3 个高信号事件。}
\vspace{3mm}
{\Large\bfseries 在职课程培训班}\par
首页 学院概况 学院简介 数说学院 现任领导 学术委员会 行政机构 新闻公告 学院新闻 学院公告 师生风采 光影视界 师资队伍 专任教师 博士后队伍 人才招聘 人才培养 本科生 学术型硕士研究生 专业学位硕士研究生 学术型博士研究生 专业学位博士研究生 在职课程培训班 同等学力申请硕士学位 科学研究 研究机构 研究成果 合作交流 对外合作 国际交流 校友名录 光影记录 校友名录 高礼研究院 EN 教师入口 人才培养 本科生 学术型硕士研究生 专业学位硕士研究生 学术型博士研究生 专业学位博士研究生 在职课程培训班 同等学力申请硕士学位 在职课程培训班 您当前所在的位置： 首页 > 人才培养 > 在职课程培训班 19 2026-10 源自人大高瓴，成就AI英才 2026年（秋）人工智能在职课程培训班招生简章 19 2026-10 人大高瓴人工智能学院人工智能专业在职课程培训的10个常见咨询问答 01 2026-03 2026年中国
\textbf{与你有关：} 这是当前用户可实际参与或申请的高可触达机会；参与范围：北京 / 人大；命中用户订阅版面：opportunity、research、ai；来源接近一手且权威度高；标题/正文包含高信号变化
\textbf{建议：} 检查报名资格、截止时间、参与形式和预期收益；合适则立即加入任务系统。
{\footnotesize\color{ormuted} 来源：中国人民大学高瓴人工智能学院公告 · \url{https://ai.ruc.edu.cn/rcpy/zzkcpxb/index.htm}}
\vspace{4mm}
\hrule height 0.25pt
\vspace{3mm}
\begin{multicols}{2}
{\large\bfseries 姚班培养-清华大学交叉信息研究院}\par
首页 热点新闻 讲座预告 院系概况 院系简介 院长寄语 现任领导 院系沿革 新闻动态 院系动态 媒体聚焦 人员队伍 全职教师 研究系列 荣誉客座 兼职教师 博士后 行政人员 实验员 姚班本科生 姚班培养 姚班招生 研究生 研究生培养 研究生招生 量子中心 中心简介 师资力量 科研概况 量子中心动态 科学研究 组织架构 课题组介绍 研究动态 科研简报 校园生活 学生活动 衣食住行 公开招聘 教师招聘 博士后招聘 职员招聘 校友频道 校友风采 校友俱乐部 校友捐赠 姚班本科生 姚班培养 姚班招生 清华学堂计算机科学实验班 姚班介绍 “清华学堂计算机科学实验班（姚班）”由世界著名计算机科学家姚期智院士于2005年创办，是清华大学培养国际顶尖计算机科学及交叉创新人才的本科生实验班，被誉为“最优秀的本科生和最优秀的本科教育”标杆。 姚班面向科技发展前沿与国家重大需求构建培养体系，以计算机科学、人工智能和量子信息三大前沿专业方向为核心，供学
\textbf{与你有关：} 这是当前用户可实际参与或申请的高可触达机会；参与范围：北京 / 清华；命中用户订阅版面：opportunity、research、ai；来源接近一手且权威度高；标题/正文包含高信号变化
\textbf{建议：} 检查报名资格、截止时间、参与形式和预期收益；合适则立即加入任务系统。
{\footnotesize\color{ormuted} 来源：清华交叉信息研究院 · \url{https://iiis.tsinghua.edu.cn/ybbks/ybpy.htm}}
\vspace{2mm}
\hrule height 0.2pt
\vspace{2mm}
{\large\bfseries Back to Careers}\par
Back to Careers Join our team All Teams AI Research \& Engineering Applied AI Communications Compute Data Science \& Analytics Engineering \& Design - Product Finance Hardware Legal Marketing \& Brand People Product Management, Support, \& Operations Public Benefit Public Policy Safeguards (Trust \& Safety) Sales Security Software Engineering - Infrastructure Technical Education Technical Program Management All Offices Ban
\textbf{与你有关：} 参与范围：全球 / 多为海外全职；命中用户订阅版面：opportunity、ai、research；来源接近一手且权威度高；标题/正文包含高信号变化
\textbf{建议：} 检查报名资格、截止时间、参与形式和预期收益；合适则立即加入任务系统。
{\footnotesize\color{ormuted} 来源：Anthropic Careers · \url{https://www.anthropic.com/careers}}
\end{multicols}
\newpage
{\small\bfseries OR MORNING / SECTION 03 / 2026-10-19}
\vspace{2mm}
\hrule height 0.9pt
\vspace{1.5mm}
\hrule height 0.25pt
\vspace{4mm}
\section*{科研前沿}
{\color{ormuted} 真正会改变你接下来怎么做研究、怎么写、怎么投的信号。}
\vspace{2mm}
{\LARGE\bfseries 在职课程培训班}\par
\vspace{1mm}
{\color{ormuted} 本版筛出 5 个高信号事件。}
\vspace{3mm}
{\Large\bfseries 在职课程培训班}\par
首页 学院概况 学院简介 数说学院 现任领导 学术委员会 行政机构 新闻公告 学院新闻 学院公告 师生风采 光影视界 师资队伍 专任教师 博士后队伍 人才招聘 人才培养 本科生 学术型硕士研究生 专业学位硕士研究生 学术型博士研究生 专业学位博士研究生 在职课程培训班 同等学力申请硕士学位 科学研究 研究机构 研究成果 合作交流 对外合作 国际交流 校友名录 光影记录 校友名录 高礼研究院 EN 教师入口 人才培养 本科生 学术型硕士研究生 专业学位硕士研究生 学术型博士研究生 专业学位博士研究生 在职课程培训班 同等学力申请硕士学位 在职课程培训班 您当前所在的位置： 首页 > 人才培养 > 在职课程培训班 19 2026-10 源自人大高瓴，成就AI英才 2026年（秋）人工智能在职课程培训班招生简章 19 2026-10 人大高瓴人工智能学院人工智能专业在职课程培训的10个常见咨询问答 01 2026-03 2026年中国
\textbf{与你有关：} 这是当前用户可实际参与或申请的高可触达机会；参与范围：北京 / 人大；命中用户订阅版面：opportunity、research、ai；来源接近一手且权威度高；标题/正文包含高信号变化
\textbf{建议：} 检查报名资格、截止时间、参与形式和预期收益；合适则立即加入任务系统。
{\footnotesize\color{ormuted} 来源：中国人民大学高瓴人工智能学院公告 · \url{https://ai.ruc.edu.cn/rcpy/zzkcpxb/index.htm}}
\vspace{4mm}
\hrule height 0.25pt
\vspace{3mm}
\begin{multicols}{2}
{\large\bfseries 姚班培养-清华大学交叉信息研究院}\par
首页 热点新闻 讲座预告 院系概况 院系简介 院长寄语 现任领导 院系沿革 新闻动态 院系动态 媒体聚焦 人员队伍 全职教师 研究系列 荣誉客座 兼职教师 博士后 行政人员 实验员 姚班本科生 姚班培养 姚班招生 研究生 研究生培养 研究生招生 量子中心 中心简介 师资力量 科研概况 量子中心动态 科学研究 组织架构 课题组介绍 研究动态 科研简报 校园生活 学生活动 衣食住行 公开招聘 教师招聘 博士后招聘 职员招聘 校友频道 校友风采 校友俱乐部 校友捐赠 姚班本科生 姚班培养 姚班招生 清华学堂计算机科学实验班 姚班介绍 “清华学堂计算机科学实验班（姚班）”由世界著名计算机科学家姚期智院士于2005年创办，是清华大学培养国际顶尖计算机科学及交叉创新人才的本科生实验班，被誉为“最优秀的本科生和最优秀的本科教育”标杆。 姚班面向科技发展前沿与国家重大需求构建培养体系，以计算机科学、人工智能和量子信息三大前沿专业方向为核心，供学
\textbf{与你有关：} 这是当前用户可实际参与或申请的高可触达机会；参与范围：北京 / 清华；命中用户订阅版面：opportunity、research、ai；来源接近一手且权威度高；标题/正文包含高信号变化
\textbf{建议：} 检查报名资格、截止时间、参与形式和预期收益；合适则立即加入任务系统。
{\footnotesize\color{ormuted} 来源：清华交叉信息研究院 · \url{https://iiis.tsinghua.edu.cn/ybbks/ybpy.htm}}
\vspace{2mm}
\hrule height 0.2pt
\vspace{2mm}
{\large\bfseries Back to Careers}\par
Back to Careers Join our team All Teams AI Research \& Engineering Applied AI Communications Compute Data Science \& Analytics Engineering \& Design - Product Finance Hardware Legal Marketing \& Brand People Product Management, Support, \& Operations Public Benefit Public Policy Safeguards (Trust \& Safety) Sales Security Software Engineering - Infrastructure Technical Education Technical Program Management All Offices Ban
\textbf{与你有关：} 参与范围：全球 / 多为海外全职；命中用户订阅版面：opportunity、ai、research；来源接近一手且权威度高；标题/正文包含高信号变化
\textbf{建议：} 检查报名资格、截止时间、参与形式和预期收益；合适则立即加入任务系统。
{\footnotesize\color{ormuted} 来源：Anthropic Careers · \url{https://www.anthropic.com/careers}}
\vspace{2mm}
\hrule height 0.2pt
\vspace{2mm}
{\large\bfseries Newsroom \textbackslash{} Anthropic}\par
Newsroom Press inquiries press@anthropic.com Non-media inquiries How to get support Media assets Download press kit Expanding the Cyber Verification Program Announcements Oct 6, 2026 We’re launching a new, expanded version of our Cyber Verification Program, which makes advanced cyber capabilities and reduced blocking classifiers available to qualifying security professionals. Announcements Oct 7, 2026 Introducing Cla
\textbf{与你有关：} 命中用户订阅版面：ai、research、math；来源接近一手且权威度高；标题/正文包含高信号变化
{\footnotesize\color{ormuted} 来源：Anthropic Research · \url{https://www.anthropic.com/news}}
\vspace{2mm}
\hrule height 0.2pt
\vspace{2mm}
{\large\bfseries Algorithms \& theory}\par
Algorithms and theory Google’s mission presents many exciting algorithmic and optimization challenges across different product areas including Search, Ads, Social, and Google Infrastructure. These include optimizing internal systems such as scheduling the machines that power the numerous computations done each day, as well as optimizations that affect core products and users, from online allocation of ads to page-vie
\textbf{与你有关：} 命中用户订阅版面：ai、research、math；来源接近一手且权威度高；标题/正文包含高信号变化
{\footnotesize\color{ormuted} 来源：Google Research Blog · \url{https://research.google/research-areas/algorithms-and-theory}}
\end{multicols}
\newpage
{\small\bfseries OR MORNING / SECTION 04 / 2026-10-19}
\vspace{2mm}
\hrule height 0.9pt
\vspace{1.5mm}
\hrule height 0.25pt
\vspace{4mm}
\section*{AI Frontier}
{\color{ormuted} Frontier model、Agent、RL、Reasoning、Safety。}
\vspace{2mm}
{\LARGE\bfseries 在职课程培训班}\par
\vspace{1mm}
{\color{ormuted} 本版筛出 5 个高信号事件。}
\vspace{3mm}
{\Large\bfseries 在职课程培训班}\par
首页 学院概况 学院简介 数说学院 现任领导 学术委员会 行政机构 新闻公告 学院新闻 学院公告 师生风采 光影视界 师资队伍 专任教师 博士后队伍 人才招聘 人才培养 本科生 学术型硕士研究生 专业学位硕士研究生 学术型博士研究生 专业学位博士研究生 在职课程培训班 同等学力申请硕士学位 科学研究 研究机构 研究成果 合作交流 对外合作 国际交流 校友名录 光影记录 校友名录 高礼研究院 EN 教师入口 人才培养 本科生 学术型硕士研究生 专业学位硕士研究生 学术型博士研究生 专业学位博士研究生 在职课程培训班 同等学力申请硕士学位 在职课程培训班 您当前所在的位置： 首页 > 人才培养 > 在职课程培训班 19 2026-10 源自人大高瓴，成就AI英才 2026年（秋）人工智能在职课程培训班招生简章 19 2026-10 人大高瓴人工智能学院人工智能专业在职课程培训的10个常见咨询问答 01 2026-03 2026年中国
\textbf{与你有关：} 这是当前用户可实际参与或申请的高可触达机会；参与范围：北京 / 人大；命中用户订阅版面：opportunity、research、ai；来源接近一手且权威度高；标题/正文包含高信号变化
\textbf{建议：} 检查报名资格、截止时间、参与形式和预期收益；合适则立即加入任务系统。
{\footnotesize\color{ormuted} 来源：中国人民大学高瓴人工智能学院公告 · \url{https://ai.ruc.edu.cn/rcpy/zzkcpxb/index.htm}}
\vspace{4mm}
\hrule height 0.25pt
\vspace{3mm}
\begin{multicols}{2}
{\large\bfseries 姚班培养-清华大学交叉信息研究院}\par
首页 热点新闻 讲座预告 院系概况 院系简介 院长寄语 现任领导 院系沿革 新闻动态 院系动态 媒体聚焦 人员队伍 全职教师 研究系列 荣誉客座 兼职教师 博士后 行政人员 实验员 姚班本科生 姚班培养 姚班招生 研究生 研究生培养 研究生招生 量子中心 中心简介 师资力量 科研概况 量子中心动态 科学研究 组织架构 课题组介绍 研究动态 科研简报 校园生活 学生活动 衣食住行 公开招聘 教师招聘 博士后招聘 职员招聘 校友频道 校友风采 校友俱乐部 校友捐赠 姚班本科生 姚班培养 姚班招生 清华学堂计算机科学实验班 姚班介绍 “清华学堂计算机科学实验班（姚班）”由世界著名计算机科学家姚期智院士于2005年创办，是清华大学培养国际顶尖计算机科学及交叉创新人才的本科生实验班，被誉为“最优秀的本科生和最优秀的本科教育”标杆。 姚班面向科技发展前沿与国家重大需求构建培养体系，以计算机科学、人工智能和量子信息三大前沿专业方向为核心，供学
\textbf{与你有关：} 这是当前用户可实际参与或申请的高可触达机会；参与范围：北京 / 清华；命中用户订阅版面：opportunity、research、ai；来源接近一手且权威度高；标题/正文包含高信号变化
\textbf{建议：} 检查报名资格、截止时间、参与形式和预期收益；合适则立即加入任务系统。
{\footnotesize\color{ormuted} 来源：清华交叉信息研究院 · \url{https://iiis.tsinghua.edu.cn/ybbks/ybpy.htm}}
\vspace{2mm}
\hrule height 0.2pt
\vspace{2mm}
{\large\bfseries Back to Careers}\par
Back to Careers Join our team All Teams AI Research \& Engineering Applied AI Communications Compute Data Science \& Analytics Engineering \& Design - Product Finance Hardware Legal Marketing \& Brand People Product Management, Support, \& Operations Public Benefit Public Policy Safeguards (Trust \& Safety) Sales Security Software Engineering - Infrastructure Technical Education Technical Program Management All Offices Ban
\textbf{与你有关：} 参与范围：全球 / 多为海外全职；命中用户订阅版面：opportunity、ai、research；来源接近一手且权威度高；标题/正文包含高信号变化
\textbf{建议：} 检查报名资格、截止时间、参与形式和预期收益；合适则立即加入任务系统。
{\footnotesize\color{ormuted} 来源：Anthropic Careers · \url{https://www.anthropic.com/careers}}
\vspace{2mm}
\hrule height 0.2pt
\vspace{2mm}
{\large\bfseries Newsroom \textbackslash{} Anthropic}\par
Newsroom Press inquiries press@anthropic.com Non-media inquiries How to get support Media assets Download press kit Expanding the Cyber Verification Program Announcements Oct 6, 2026 We’re launching a new, expanded version of our Cyber Verification Program, which makes advanced cyber capabilities and reduced blocking classifiers available to qualifying security professionals. Announcements Oct 7, 2026 Introducing Cla
\textbf{与你有关：} 命中用户订阅版面：ai、research、math；来源接近一手且权威度高；标题/正文包含高信号变化
{\footnotesize\color{ormuted} 来源：Anthropic Research · \url{https://www.anthropic.com/news}}
\vspace{2mm}
\hrule height 0.2pt
\vspace{2mm}
{\large\bfseries Algorithms \& theory}\par
Algorithms and theory Google’s mission presents many exciting algorithmic and optimization challenges across different product areas including Search, Ads, Social, and Google Infrastructure. These include optimizing internal systems such as scheduling the machines that power the numerous computations done each day, as well as optimizations that affect core products and users, from online allocation of ads to page-vie
\textbf{与你有关：} 命中用户订阅版面：ai、research、math；来源接近一手且权威度高；标题/正文包含高信号变化
{\footnotesize\color{ormuted} 来源：Google Research Blog · \url{https://research.google/research-areas/algorithms-and-theory}}
\end{multicols}
\newpage
{\small\bfseries OR MORNING / SECTION 05 / 2026-10-19}
\vspace{2mm}
\hrule height 0.9pt
\vspace{1.5mm}
\hrule height 0.25pt
\vspace{4mm}
\section*{数学 × AI}
{\color{ormuted} AI for Mathematics、Lean、证明与重要数学进展。}
\vspace{2mm}
{\LARGE\bfseries Newsroom \textbackslash{} Anthropic}\par
\vspace{1mm}
{\color{ormuted} 本版筛出 5 个高信号事件。}
\vspace{3mm}
{\Large\bfseries Newsroom \textbackslash{} Anthropic}\par
Newsroom Press inquiries press@anthropic.com Non-media inquiries How to get support Media assets Download press kit Expanding the Cyber Verification Program Announcements Oct 6, 2026 We’re launching a new, expanded version of our Cyber Verification Program, which makes advanced cyber capabilities and reduced blocking classifiers available to qualifying security professionals. Announcements Oct 7, 2026 Introducing Cla
\textbf{与你有关：} 命中用户订阅版面：ai、research、math；来源接近一手且权威度高；标题/正文包含高信号变化
{\footnotesize\color{ormuted} 来源：Anthropic Research · \url{https://www.anthropic.com/news}}
\vspace{4mm}
\hrule height 0.25pt
\vspace{3mm}
\begin{multicols}{2}
{\large\bfseries Algorithms \& theory}\par
Algorithms and theory Google’s mission presents many exciting algorithmic and optimization challenges across different product areas including Search, Ads, Social, and Google Infrastructure. These include optimizing internal systems such as scheduling the machines that power the numerous computations done each day, as well as optimizations that affect core products and users, from online allocation of ads to page-vie
\textbf{与你有关：} 命中用户订阅版面：ai、research、math；来源接近一手且权威度高；标题/正文包含高信号变化
{\footnotesize\color{ormuted} 来源：Google Research Blog · \url{https://research.google/research-areas/algorithms-and-theory}}
\vspace{2mm}
\hrule height 0.2pt
\vspace{2mm}
{\large\bfseries v4.36.0-rc1}\par
v4.36.0-rc1 Pre-release Pre-release Compare Choose a tag to compare Sorry, something went wrong. Filter Loading Sorry, something went wrong. Uh oh! There was an error while loading. Please reload this page . No results found View all tags github-actions released this 16 Oct 19:43 · 77 commits to master since this release v4.36.0-rc1 625d891 chore: bump toolchain to v4.36.0-rc1 (\#44879) Co-authored-by: Anne C.A. Baane
\textbf{与你有关：} 命中用户订阅版面：math、engineering、research；来源接近一手且权威度高；标题/正文包含高信号变化
{\footnotesize\color{ormuted} 来源：mathlib Releases · \url{https://github.com/leanprover-community/mathlib4/releases/tag/v4.36.0-rc1}}
\vspace{2mm}
\hrule height 0.2pt
\vspace{2mm}
{\large\bfseries Proof-based general-purpose computing with Magnetar}\par
At a glance A new form of computing . Zero-knowledge proofs let you run a workload and receive, alongside the result, a short proof that the program ran correctly. Anyone can check the proof in milliseconds, without rerunning the program, without trusting whoever ran it, and without needing data on which the program was run. Three grand challenges have stood in the way of broad adoption : making proofs fast enough to
\textbf{与你有关：} 命中用户订阅版面：ai、research、engineering；来源接近一手且权威度高；标题/正文包含高信号变化
{\footnotesize\color{ormuted} 来源：Microsoft Research Blog · \url{https://www.microsoft.com/en-us/research/blog/proof-based-general-purpose-computing-with-magnetar/}}
\vspace{2mm}
\hrule height 0.2pt
\vspace{2mm}
{\large\bfseries Scaling Indic Pretraining: 2.8 Trillion Tokens and an Open Data Generation Pipeline}\par
Back to Articles Scaling Indic Pretraining: 2.8 Trillion Tokens and an Open Data Generation Pipeline Enterprise + Article Published October 16, 2026 Upvote 32 +26 Will Jennings WillJenningsDC Follow nvidia Kiran Praveen kipraveen Follow nvidia Sriharsha Niverty sniverty Follow nvidia Raviraj Joshi ravirajoshi Follow nvidia Markus Kliegl mkliegl-nv Follow nvidia Yev Meyer nv-3mei Follow nvidia Utkarsh Vaidya uvaidya F
\textbf{与你有关：} 命中用户订阅版面：engineering、ai、research；来源接近一手且权威度高；标题/正文包含高信号变化
{\footnotesize\color{ormuted} 来源：Hugging Face Blog · \url{https://huggingface.co/blog/nvidia/scaling-indic-pretraining}}
\end{multicols}
\newpage
{\small\bfseries OR MORNING / SECTION 06 / 2026-10-19}
\vspace{2mm}
\hrule height 0.9pt
\vspace{1.5mm}
\hrule height 0.25pt
\vspace{4mm}
\section*{Engineering}
{\color{ormuted} 开源、框架、GPU、推理基础设施与真正值得使用的新工具。}
\vspace{2mm}
{\LARGE\bfseries Back to Careers}\par
\vspace{1mm}
{\color{ormuted} 本版筛出 5 个高信号事件。}
\vspace{3mm}
{\Large\bfseries Back to Careers}\par
Back to Careers Join our team All Teams AI Research \& Engineering Applied AI Communications Compute Data Science \& Analytics Engineering \& Design - Product Finance Hardware Legal Marketing \& Brand People Product Management, Support, \& Operations Public Benefit Public Policy Safeguards (Trust \& Safety) Sales Security Software Engineering - Infrastructure Technical Education Technical Program Management All Offices Ban
\textbf{与你有关：} 参与范围：全球 / 多为海外全职；命中用户订阅版面：opportunity、ai、research；来源接近一手且权威度高；标题/正文包含高信号变化
\textbf{建议：} 检查报名资格、截止时间、参与形式和预期收益；合适则立即加入任务系统。
{\footnotesize\color{ormuted} 来源：Anthropic Careers · \url{https://www.anthropic.com/careers}}
\vspace{4mm}
\hrule height 0.25pt
\vspace{3mm}
\begin{multicols}{2}
{\large\bfseries Algorithms \& theory}\par
Algorithms and theory Google’s mission presents many exciting algorithmic and optimization challenges across different product areas including Search, Ads, Social, and Google Infrastructure. These include optimizing internal systems such as scheduling the machines that power the numerous computations done each day, as well as optimizations that affect core products and users, from online allocation of ads to page-vie
\textbf{与你有关：} 命中用户订阅版面：ai、research、math；来源接近一手且权威度高；标题/正文包含高信号变化
{\footnotesize\color{ormuted} 来源：Google Research Blog · \url{https://research.google/research-areas/algorithms-and-theory}}
\vspace{2mm}
\hrule height 0.2pt
\vspace{2mm}
{\large\bfseries v4.36.0-rc1}\par
v4.36.0-rc1 Pre-release Pre-release Compare Choose a tag to compare Sorry, something went wrong. Filter Loading Sorry, something went wrong. Uh oh! There was an error while loading. Please reload this page . No results found View all tags github-actions released this 16 Oct 19:43 · 77 commits to master since this release v4.36.0-rc1 625d891 chore: bump toolchain to v4.36.0-rc1 (\#44879) Co-authored-by: Anne C.A. Baane
\textbf{与你有关：} 命中用户订阅版面：math、engineering、research；来源接近一手且权威度高；标题/正文包含高信号变化
{\footnotesize\color{ormuted} 来源：mathlib Releases · \url{https://github.com/leanprover-community/mathlib4/releases/tag/v4.36.0-rc1}}
\vspace{2mm}
\hrule height 0.2pt
\vspace{2mm}
{\large\bfseries Offloaded inference for real-world physical AI robotics}\par
At a glance Challenges a core assumption in robotics AI: Our research shows that running physical AI inference exclusively on onboard GPUs can limit robot performance, battery life, and scalability, and that offloading inference to edge or cloud GPUs can offer significant advantages. Demonstrates measurable benefits of inference offloading: Across representative mobile manipulation workloads, offloading improved task
\textbf{与你有关：} 命中用户订阅版面：ai、research、engineering；来源接近一手且权威度高；标题/正文包含高信号变化
{\footnotesize\color{ormuted} 来源：Microsoft Research Blog · \url{https://www.microsoft.com/en-us/research/blog/offloaded-inference-for-real-world-physical-ai-robotics/}}
\vspace{2mm}
\hrule height 0.2pt
\vspace{2mm}
{\large\bfseries Agent Lightning v1.0: A 3,500-Line Lightweight Agentic RL Framework for Training Agents with Real Harnesses}\par
At a glance Harnessed Agentic RL: Microsoft Research Asia introduces a training paradigm in which the same agent harness used in deployment participates directly in reinforcement learning, removing the need to reimplement the agent inside the training framework. Lightweight by design: Agent Lightning v1.0 delivers a complete agent RL control plane in roughly 3,500 lines of code. Native Kubernetes support: agents run 
\textbf{与你有关：} 命中用户订阅版面：ai、research、engineering；来源接近一手且权威度高；标题/正文包含高信号变化
{\footnotesize\color{ormuted} 来源：Microsoft Research Blog · \url{https://www.microsoft.com/en-us/research/blog/agent-lightning-v1-0-a-3500-line-lightweight-agentic-rl-framework-for-training-agents-with-real-harnesses/}}
\end{multicols}
\newpage
{\small\bfseries OR MORNING / SECTION 07 / 2026-10-19}
\vspace{2mm}
\hrule height 0.9pt
\vspace{1.5mm}
\hrule height 0.25pt
\vspace{4mm}
\section*{AI 商业 / 产业}
{\color{ormuted} 公司、融资、市场与可能改变 AI 发展路径的产业信号。}
\vspace{2mm}
{\LARGE\bfseries Back to Careers}\par
\vspace{1mm}
{\color{ormuted} 本版筛出 5 个高信号事件。}
\vspace{3mm}
{\Large\bfseries Back to Careers}\par
Back to Careers Join our team All Teams AI Research \& Engineering Applied AI Communications Compute Data Science \& Analytics Engineering \& Design - Product Finance Hardware Legal Marketing \& Brand People Product Management, Support, \& Operations Public Benefit Public Policy Safeguards (Trust \& Safety) Sales Security Software Engineering - Infrastructure Technical Education Technical Program Management All Offices Ban
\textbf{与你有关：} 参与范围：全球 / 多为海外全职；命中用户订阅版面：opportunity、ai、research；来源接近一手且权威度高；标题/正文包含高信号变化
\textbf{建议：} 检查报名资格、截止时间、参与形式和预期收益；合适则立即加入任务系统。
{\footnotesize\color{ormuted} 来源：Anthropic Careers · \url{https://www.anthropic.com/careers}}
\vspace{4mm}
\hrule height 0.25pt
\vspace{3mm}
\begin{multicols}{2}
{\large\bfseries Scaling Indic Pretraining: 2.8 Trillion Tokens and an Open Data Generation Pipeline}\par
Back to Articles Scaling Indic Pretraining: 2.8 Trillion Tokens and an Open Data Generation Pipeline Enterprise + Article Published October 16, 2026 Upvote 32 +26 Will Jennings WillJenningsDC Follow nvidia Kiran Praveen kipraveen Follow nvidia Sriharsha Niverty sniverty Follow nvidia Raviraj Joshi ravirajoshi Follow nvidia Markus Kliegl mkliegl-nv Follow nvidia Yev Meyer nv-3mei Follow nvidia Utkarsh Vaidya uvaidya F
\textbf{与你有关：} 命中用户订阅版面：engineering、ai、research；来源接近一手且权威度高；标题/正文包含高信号变化
{\footnotesize\color{ormuted} 来源：Hugging Face Blog · \url{https://huggingface.co/blog/nvidia/scaling-indic-pretraining}}
\vspace{2mm}
\hrule height 0.2pt
\vspace{2mm}
{\large\bfseries Self-Building Labs}\par
1]Department of Chemistry, University of Toronto, 80 St. George St., Toronto, ON M5S 3H6, Canada 2]Department of Computer Science, University of Toronto, 40 St George St., Toronto, ON M5S 2E4, Canada 3]Vector Institute for Artificial Intelligence, W1140-108 College St., Schwartz Reisman Innovation Campus, Toronto, ON M5G 0C6, Canada 4]Department of Materials Science \& Engineering, University of Toronto, 184 College S
\textbf{与你有关：} 命中用户订阅版面：ai、research、math；来源接近一手且权威度高；标题/正文包含高信号变化
{\footnotesize\color{ormuted} 来源：arXiv cs.AI Recent · \url{https://arxiv.org/html/2610.23826v1}}
\vspace{2mm}
\hrule height 0.2pt
\vspace{2mm}
{\large\bfseries arXiv:2610.23796}\par
Computer Science > Artificial Intelligence arXiv:2610.23796 (cs) [Submitted on 16 Oct 2026] Title: FuTransformer: FP8 prefill and INT4 decode for Llama-style inference on Hopper Authors: Wei Wang View a PDF of the paper titled FuTransformer: FP8 prefill and INT4 decode for Llama-style inference on Hopper, by Wei Wang View PDF HTML (experimental) Abstract: Transformer prefill reuses weights across tokens and is domina
\textbf{与你有关：} 命中用户订阅版面：ai、research、math；来源接近一手且权威度高；标题/正文包含高信号变化
{\footnotesize\color{ormuted} 来源：arXiv cs.AI Recent · \url{https://arxiv.org/abs/2610.23796}}
\vspace{2mm}
\hrule height 0.2pt
\vspace{2mm}
{\large\bfseries 1Introduction}\par
marginparsep has been altered. topmargin has been altered. marginparwidth has been altered. marginparpush has been altered. The page layout violates the ICML style. Please do not change the page layout, or include packages like geometry, savetrees, or fullpage, which change it for you. We’re not able to reliably undo arbitrary changes to the style. Please remove the offending package(s), or layout-changing commands a
\textbf{与你有关：} 命中用户订阅版面：ai、research、math；来源接近一手且权威度高；标题/正文包含高信号变化
{\footnotesize\color{ormuted} 来源：arXiv cs.AI Recent · \url{https://arxiv.org/html/2610.23796v1}}
\end{multicols}
\newpage
{\small\bfseries OR MORNING / SECTION 08 / 2026-10-19}
\vspace{2mm}
\hrule height 0.9pt
\vspace{1.5mm}
\hrule height 0.25pt
\vspace{4mm}
\section*{世界状态}
{\color{ormuted} 重大政策、国际、宏观与不应因为沉浸在 AI 中而错过的变化。}
\vspace{2mm}
{\LARGE\bfseries 在职课程培训班}\par
\vspace{1mm}
{\color{ormuted} 本版筛出 3 个高信号事件。}
\vspace{3mm}
{\Large\bfseries 在职课程培训班}\par
首页 学院概况 学院简介 数说学院 现任领导 学术委员会 行政机构 新闻公告 学院新闻 学院公告 师生风采 光影视界 师资队伍 专任教师 博士后队伍 人才招聘 人才培养 本科生 学术型硕士研究生 专业学位硕士研究生 学术型博士研究生 专业学位博士研究生 在职课程培训班 同等学力申请硕士学位 科学研究 研究机构 研究成果 合作交流 对外合作 国际交流 校友名录 光影记录 校友名录 高礼研究院 EN 教师入口 人才培养 本科生 学术型硕士研究生 专业学位硕士研究生 学术型博士研究生 专业学位博士研究生 在职课程培训班 同等学力申请硕士学位 在职课程培训班 您当前所在的位置： 首页 > 人才培养 > 在职课程培训班 19 2026-10 源自人大高瓴，成就AI英才 2026年（秋）人工智能在职课程培训班招生简章 19 2026-10 人大高瓴人工智能学院人工智能专业在职课程培训的10个常见咨询问答 01 2026-03 2026年中国
\textbf{与你有关：} 这是当前用户可实际参与或申请的高可触达机会；参与范围：北京 / 人大；命中用户订阅版面：opportunity、research、ai；来源接近一手且权威度高；标题/正文包含高信号变化
\textbf{建议：} 检查报名资格、截止时间、参与形式和预期收益；合适则立即加入任务系统。
{\footnotesize\color{ormuted} 来源：中国人民大学高瓴人工智能学院公告 · \url{https://ai.ruc.edu.cn/rcpy/zzkcpxb/index.htm}}
\vspace{4mm}
\hrule height 0.25pt
\vspace{3mm}
\begin{multicols}{2}
{\large\bfseries 姚班培养-清华大学交叉信息研究院}\par
首页 热点新闻 讲座预告 院系概况 院系简介 院长寄语 现任领导 院系沿革 新闻动态 院系动态 媒体聚焦 人员队伍 全职教师 研究系列 荣誉客座 兼职教师 博士后 行政人员 实验员 姚班本科生 姚班培养 姚班招生 研究生 研究生培养 研究生招生 量子中心 中心简介 师资力量 科研概况 量子中心动态 科学研究 组织架构 课题组介绍 研究动态 科研简报 校园生活 学生活动 衣食住行 公开招聘 教师招聘 博士后招聘 职员招聘 校友频道 校友风采 校友俱乐部 校友捐赠 姚班本科生 姚班培养 姚班招生 清华学堂计算机科学实验班 姚班介绍 “清华学堂计算机科学实验班（姚班）”由世界著名计算机科学家姚期智院士于2005年创办，是清华大学培养国际顶尖计算机科学及交叉创新人才的本科生实验班，被誉为“最优秀的本科生和最优秀的本科教育”标杆。 姚班面向科技发展前沿与国家重大需求构建培养体系，以计算机科学、人工智能和量子信息三大前沿专业方向为核心，供学
\textbf{与你有关：} 这是当前用户可实际参与或申请的高可触达机会；参与范围：北京 / 清华；命中用户订阅版面：opportunity、research、ai；来源接近一手且权威度高；标题/正文包含高信号变化
\textbf{建议：} 检查报名资格、截止时间、参与形式和预期收益；合适则立即加入任务系统。
{\footnotesize\color{ormuted} 来源：清华交叉信息研究院 · \url{https://iiis.tsinghua.edu.cn/ybbks/ybpy.htm}}
\vspace{2mm}
\hrule height 0.2pt
\vspace{2mm}
{\large\bfseries Back to Careers}\par
Back to Careers Join our team All Teams AI Research \& Engineering Applied AI Communications Compute Data Science \& Analytics Engineering \& Design - Product Finance Hardware Legal Marketing \& Brand People Product Management, Support, \& Operations Public Benefit Public Policy Safeguards (Trust \& Safety) Sales Security Software Engineering - Infrastructure Technical Education Technical Program Management All Offices Ban
\textbf{与你有关：} 参与范围：全球 / 多为海外全职；命中用户订阅版面：opportunity、ai、research；来源接近一手且权威度高；标题/正文包含高信号变化
\textbf{建议：} 检查报名资格、截止时间、参与形式和预期收益；合适则立即加入任务系统。
{\footnotesize\color{ormuted} 来源：Anthropic Careers · \url{https://www.anthropic.com/careers}}
\end{multicols}
\end{document}
